\section{Problem Formulation}

In the considered OIRS-assisted VLC system, the objective is to determine the
orientation of the optical intelligent reflecting surface (OIRS) that maximizes
the achievable data rate (ADR). The OIRS orientation is characterized by the
rotation angles $\gamma$ and $\omega$, which directly determine the effective
optical channel gain $z(\gamma,\omega)$ and, consequently, the communication
performance.

Based on the capacity lower bound for intensity-modulation/direct-detection
(IM/DD) channels, the ADR is expressed as
\begin{equation}
\mathrm{ADR}(\gamma,\omega)
= \frac{B}{2}\log_{2}\!\left[
1+\frac{e}{2\pi}
\left(\frac{P\,R_{\mathrm{PD}}\,z(\gamma,\omega)}{Q}\right)^{2}
\right],
\label{eq:adr}
\end{equation}
where $B$ denotes the modulation bandwidth, $P$ is the transmitted optical
power, $R_{\mathrm{PD}}$ is the photodetector responsivity, $Q$ is the noise
standard deviation at the receiver, and $z(\gamma,\omega)$ is the effective
channel gain determined by the OIRS orientation.

Accordingly, the OIRS orientation design problem is formulated as
\begin{equation}
\begin{aligned}
(\mathrm{P1}):\quad
(\gamma^{*},\omega^{*})
= \arg\max_{\gamma,\,\omega}\;
& \mathrm{ADR}(\gamma,\omega) \\
\text{s.t.}\quad
& \gamma_{\min} \le \gamma \le \gamma_{\max}, \\
& \omega_{\min} \le \omega \le \omega_{\max}.
\end{aligned}
\label{eq:p1}
\end{equation}

Since the optimization variables $\gamma$ and $\omega$ are continuous and the
mapping from the OIRS orientation to the channel gain is highly nonlinear and
non-convex, obtaining a closed-form solution to (P1) is intractable. Therefore,
(P1) is solved using deep reinforcement learning (DRL), in which the optimal
orientation is learned through direct interaction with the communication
environment. At each time step, the agent observes the system state (e.g., the
transmitter, OIRS, and receiver geometry together with the current orientation
angles) and outputs the continuous action vector
\begin{equation}
\mathbf{a}=[\gamma,\ \omega].
\end{equation}
The selected orientation determines the effective channel gain
$z(\gamma,\omega)$, from which the ADR in \eqref{eq:adr} is evaluated and
returned by the environment as the reward, enabling the agent to iteratively
learn the orientation that maximizes the communication performance.

To comprehensively evaluate the proposed optimization framework, five DRL
algorithms are employed: Deep Deterministic Policy Gradient (DDPG), Soft
Actor--Critic (SAC), Proximal Policy Optimization (PPO), Continuous Action
Policy Gradient (CAPG), and the proposed Enhanced Continuous Deep Policy
Optimization (ECDPO) algorithm. All algorithms operate under an identical
system model, state representation, action space, and reward function, ensuring
a fair comparison of their convergence behavior, learning stability, and
achieved ADR. The algorithm attaining the highest ADR with stable convergence
is identified as the most suitable approach for optimal OIRS orientation
design.
